\section{Mathematical theory and no-arbitrage principle}

In this section, we will focus on the theoretical notions and concepts necessary to understand finance.

\subsection{Filtered probability space}

Consider the triplet $(\Omega,\mathcal{F},\mathbb{P})$ forming a probability space, with : 
\begin{itemize}
    \item[$\bullet$] the sample space $\Omega$ representing the possible states of a financial market
    \item[$\bullet$] a $\sigma$-algebra $\mathcal{F}$ on the sample space $\Omega$
    \item[$\bullet$] a probability measure $\mathbb{P}$ on the $\sigma$-algebra $\mathcal{F}$
\end{itemize}

Let $I = \{0,1,...,T\}$, with $T < \infty$ be the discrete set representing time points, where $T$ is the option's maturity date.
We introduce a $\textbf{filtration}$ $(\mathcal{F}_t)_{t \in I}$ on the probability space $(\Omega,\mathcal{F},\mathbb{P})$. 
The filtration $(\mathcal{F}_t)_{t \in I}$ is an increasing set of sub-$\sigma$-algebra of $\mathcal{F}$ with the property : 
\begin{align*}
    \forall (t_0,t_1) \in I \times I, \qquad t_0 \le t_1 \implies \mathcal{F}_{t_0} \subseteq \mathcal{F}_{t_1} \subseteq \mathcal{F}
\end{align*}
\cite{Filtration}
We then defin the \textbf{filtered probability space} $(\Omega, \mathcal{F}, (\mathcal{F}_t)_{t \in I}, \mathbb{P})$

\bigskip

\remark Working within this space allows us to model the evolution of a financial market.
Thus, each $\sigma$-algebra in the filtration $(\mathcal{F}_t)_{t \in I}$ represents the information known at time $t$.
Therefore, an agent at time $t_1$ has access to all information from an earlier time $t_0$, provided $t_0 \leq t_1$.
The studied stochastic process $X\coloneq(X_t)_{t \in I}$ must be adapted to the filtration $(\mathcal{F}_t)_{t \in I}$ ($\forall t \in I$, $X_t$ is $\mathcal{F}_t$-measurable) 
We can now define the random variables used to model the financial processes of the model.

Consider the following random variables :
\begin{align*}
    \text{Asset quantity : } \quad \alpha_t& : \Omega \to \mathbb{R} \\
    \text{Risky asset : } \quad S_t& : \Omega \to \mathbb{R}^+ \\
    \text{Risk-free asset : } \quad B_t& : \Omega \to \mathbb{R}^+, t \mapsto (1+r)^t
\end{align*}
with $r>-1$ the risk-free rate.
We assume that the random variables $\alpha_t$, $S_t$ and $B_t$ are $\mathcal{F}_t$-measurable for all $t \in I$.

\remark For each $t \in I$, the random variable $B_t$ is constant on $\Omega$ ($\forall \omega \in \Omega, B_t(\omega)=(1+r)^t$).
Therefore, it is a deterministic random variable.
$S_t$ represents financial market uncertainty, $\alpha_t$ is a strategy adapted to the filtration (and thus adapted to the information known at time $t$), and $B_t$ represents the deterministic evolution of a risk-free investment.



\pagebreak



\subsection{Portfolio theory}

We then introduce the concept of an \textbf{investment strategy} as a stochastic process $\Theta = (\Theta_t)_{t \in I} \coloneq (\alpha_t, \beta_t)_{t \in I}$
where $\alpha_t$ and $\beta_t$ represent respectively the number of units at time $t$ in the risky asset and in the risk-free asset.
This stochastic process is adapted to the filtration $(\mathcal{F}_t)_{t \in I}$.

Let $B_t = (1+r)^t$ be the value of one unit invested in the risk-free asset.
Thus, the value of the risk-free account is given by $\beta_t B_t$.

A \textbf{discrete portfolio $\Pi_t$} at time $t$ is defined by the value : 
\begin{align*}
    \Pi_t = \alpha_tS_t+\beta_tB_t 
\end{align*}
Therefore, the stochastic process $\Pi_t$ is composed of two parts : $\alpha_tS_t$ representing the risky component of the portfolio,
and $\beta_tB_t$ representing the risk-free component of the account. \cite{Portfolio}

Although the model evolves in discrete time, each transition from $t$ to $t+1$ contains two phases : 
\begin{itemize}
    \item Price Update : Asset prices change from $(S_t,B_t)$ to $(S_{t+1},B_{t+1})$. 
    The value of the portfolio held from period $t$, prior to any updates, becomes $\alpha_tS_{t+1}+\beta_tB_{t+1}$.
    \item Portfolio update : The investor updates their holdings from $(\alpha_t,\beta_t)$ to $(\alpha_{t+1}, \beta_{t+1})$ at the newly observed prices $(S_{t+1},B_{t+1})$.
    The value of the new portfolio is $\alpha_{t+1}S_{t+1} + \beta_{t+1}B_{t+1}$.
\end{itemize}

\bigskip

\begin{definitionbox}[frametitle = {Self-financing property}]
    A discrete portfolio is said to be \textbf{self-financing} if
    \begin{align*}
        \forall t \in I,  \quad \Pi_{t+1} = \alpha_t S_{t+1} + \beta_tB_{t+1} = \alpha_{t+1}S_{t+1} + \beta_{t+1}B_{t+1} 
    \end{align*}\cite{SelfFinancingPortfolio}
\end{definitionbox}

\bigskip

Thus, the self-financing property ensures that the evolution of $\Pi_t$ is entirely determined by the strategy $\Theta$ and the dynamics of both $S_t$ and $B_t$, without external adjustments.
In the following, we will assume that all portfolios are self-financing.

\bigskip

\begin{definitionbox}[frametitle = {No-Arbitrage Principle}]
    The No-Arbitrage principle means that there is no strategy $\Theta$ ensuring a self-financing portfolio with zero initial value and positive final value almost surely.
    \begin{center}
         $\nexists \Theta, \Pi_0=0, \; \mathbb{P}(\Pi_T\geq0)=1, \; \mathbb{P}(\Pi_T > 0) > 0$
    \end{center}
\end{definitionbox}
\bigskip



\pagebreak

\remark The No-arbitrage principle forbid the theoretical existence of a risk-free asset requiring no initial investment.
Consequently, any potential gain requires an initial outlay or entails risk.
Without this condition, every investors would exploit arbitrage opportunities, causing markets to collapse.
This property is a necessary condition to ensure market consistency.
In practice, very minor arbitrage opportunities appear during market fluctuations and are exploited by arbitrageurs. \cite{Arbitrageurs}

\subsection{Replication strategy}
\bigskip
\begin{definitionbox}[frametitle = {Replication}]
    A payoff $H$ is said to be \textbf{replicable} if : $\exists \Pi, \; \Pi_T=H \quad \mathbb{P}$-a.s. \\
    $\Pi$ is called the \textbf{replicating portfolio for H}.
\end{definitionbox}
\bigskip
\begin{definitionbox}[frametitle = {Market Completeness}]
    A market is said to be \textbf{complete} if every payoff is replicable with a portfolio.
\end{definitionbox}
\bigskip
\remark 
It follows naturally that market completeness implies the existence of a replicating portfolio for every payoff.
However, one can go further and demonstrate the uniqueness of the portolio for each payoff.

Consider a complete market, for any payoff H, there exists a \textbf{unique} replicating portfolio for H.

\begin{proof}

Let $\Pi = \Pi^1-\Pi^2$ be the difference between two self-financing portfolios that both replicate the payoff $H$.
Since the difference of two self-financing portfolios is self-financing, $\Pi$ is self-financing. Furthermore,
$$\Pi_T=\Pi_T^1-\Pi_T^2=H-H=0$$

We will now prove that $\Pi_t=0$ for every $t\in[0,T]$.

Suppose by contradiction that there exists $t\in[0,T]$ such that $\Pi_t \neq 0$.

If $\Pi_t>0$, take the opposite position $-\Pi$ at time $t$. This creates an initial cash amount $\Pi_t>0$. 
Invest this amount in the risk-free asset until time $T$. 

Since $\Pi_T=0,$ the position $-\Pi$ has zero final value, while the amount invested in the risk-free asset has strictly positive value
$$\Pi_t\frac{B_T}{B_t}>0$$

Thus, starting from zero initial wealth at time $t$, we obtain a strictly positive final payoff, yielding an arbitrage opportunity.

If $\Pi_t<0$, take the position $\Pi$ at time $t$ and borrow $-\Pi_t>0$ from the risk-free asset. The initial value of the combined position is $\Pi_t+(-\Pi_t)=0$.

At time $T$, the portfolio $\Pi$ has value $\Pi_T=0$, while the loan has to be repaid in the amount
$$(-\Pi_t)\frac{B_T}{B_t}>0$$

Therefore this position produces a strictly negative final payoff. Taking the opposite position produces an arbitrage opportunity.
In both cases, $\Pi_t \neq 0$ leads to an arbitrage opportunity, contradicting the no-arbitrage assumption.

Therefore, $\Pi_t=0\qquad\text{for all }t\in[0,T]$

Consequently, $\Pi_t^1=\Pi_t^2\qquad\text{for all }t\in[0,T]$

Thus the self-financing replicating portfolio for $H$ is unique.
\end{proof}



\subsection{Risk-neutral measure}

\bigskip

\begin{definitionbox}[frametitle = {Fundamental Theorems of Asset Pricing}]
    \textbf{First Theorem : } A discret financial market on the filtered probability space \\
    $(\Omega, \mathcal{F}, (\mathcal{F}_t)_{t \in I}, \mathbb{P})$ is arbitrage-free 
    if and only if there exists at least one risk neutral probability measure that is equivalent to the original probability measure $\mathbb{P}$.

    \textbf{Second Theorem : } An arbitrage-free market is complete if and only if there exists an unique risk-neutral measure,
    equivalent to $\mathbb{P}$ for which the \textbf{numeraire} (reference asset used to express prices) is the risk-free asset.
\end{definitionbox} \cite{FundamentalTheorems} \cite{FundamentalTheorem} \cite{Martingales}

\bigskip

\remark
We previously required markets to satisfy the no-arbitrage principle. Thus, we assume the existence of a risk-neutral measure.
Furthermore, if the market is complete, the risk-neutral measure is \textbf{unique}, allowing for the consistent pricing of all assets.
This theorem forms the basis of models for valuing options and derivatives.

The second theorem states that a market is complete if and only if there exists a unique risk-neutral probability measure.
In the binomial model, this property corresponds to the fact that the number of future states (up/down) equals the number of assets (risky/risk-free).
Thus, the model considered so far involves two types of assets (risky/risk-free), and only two possible scenarios (rise/fall) can be considered at each step.
This justifies the use of binomial models, whichis the easiest complete model, as opposed to a trionmial model, which would be incomplete.

\bigskip

\begin{definitionbox}[frametitle = {Martingale}]
    $M \coloneq (M_t)_{t \in I}$ is an $\mathcal{F}$-martingale if : 
    \begin{itemize}
        \item $\forall t \in I, M_t$ is $\mathcal{F}_t$-measurable
        \item $\forall t \in I, M_t$ is integrable 
        \item $\forall s \le t$, $M_{s} = \mathbb{E}(M_{t} | \mathcal{F}_{s})$ a.s.
    \end{itemize}
\end{definitionbox}

\bigskip

Thus, the concept of martingale is useful to represent a financial market.
Intuitively, the future expectation of a martingale, taking the past into account, is equal to the current value $(\mathbb{E}(M_{t+1}|\mathcal{F}_t) = M_t)$.
 
Consider a discret market with a risk asset $S_t$ and a risk-free asset $B_t$.
The \textbf{discounted yield} of the risky asset is defined as : 
$$\tilde{S}_t \coloneq \frac{S_t}{B_t}$$
\cite{IntroductionALevaluationDoptions}

We obviously see that the price is expressed in terms of $B_t$, which thus becomes the numeraire.
According to the first fundamental theorem of asset pricing, if the market is \textbf{arbritrage-free}, there exists at least one measure $\mathbb{Q}$ equivalent to $\mathbb{P}$ such that : 

\begin{align*}
    \tilde{S}_t = \mathbb{E}_{\mathbb{Q}}\big[\tilde{S}_{t+1} \mid \mathcal{F}_t]
\end{align*}
In other words, under $\mathbb{Q}$, the discounted price becomes an $\mathcal{F}$-martingale.

To better understand this concept, let's look at an example : 
Consider a market with : 
$$
S_0 = 100, \quad S_1 = 
\begin{cases}
120 \; \text{with probability $p$} \\
90 \; \text{with probability $1-p$} \end{cases} \quad \text{avec } r_f = 5\%
$$

We search the risk-neutral probability $q$ such that the discounted yield is a martingale under $\mathbb{Q}$:

$$\tilde{S}_0 = \frac{1}{1+r}\mathbb{E_Q}[S_1|\mathcal{F}_0] = \frac{120q + 90(1-p)}{1+r_f}$$

The solution $q$ corresponds to the \textbf{risk-neutral probability} of an asset price increase.
Under this probability, the discounted yield $\tilde{S}_t$ is a martingale and is used to value all derivative products.

  
Let $\Pi_T$ be the \textbf{payoff} of an asset at maturity $T$. 
If the market is \textbf{arbitrage-free}, the first fundamental theorem guarantees the existence of a risk-neutral measure $\mathbb{Q}$.

The initial price $\Pi_0$ of the asset is defined as the discounted expectation of its payoff under the measure $\mathbb{Q}$:
$$\Pi_0 = \mathbb{E}_{\mathbb{Q}}\Big[\frac{\Pi_T}{B_T}\Big] = \frac{\mathbb{E}_{\mathbb{Q}}[\Pi_T]}{B_T} = \frac{\mathbb{E}_{\mathbb{Q}}[\Pi_T]}{(1+r)^T}$$
This equality holds because $B$ is deterministic.

Example:
Consider a European call option with strike $K = 100$, maturity $T = 1$ period, and underlying asset $S_T$ such that:
$$
S_T = 
\begin{cases}
120 \; \text{with probability } p \\
50 \; \text{with probability } 1-p
\end{cases} \quad  \text{avec } r_f = 5\%.
$$

The option being being a call, its payoff is $\Pi_T = \max(S_T - K, 0)$.  

The initial price is :
$$ \Pi_0 = \frac{\mathbb{E}_{\mathbb{Q}}[\Pi_T]}{1+r_f} = \frac{1}{1+r_f} \big(20q + (1-q) \cdot 0 \big) = \frac{20q}{1 + r_f} $$
$q$ is the risk-neutral probability of an increase in the asset's price.
Under this probability, the discounted option price is consistent with an arbitrage-free market.